\begin{frame}\titlepage\end{frame}

\begin{frame}{Введение. Корреляционная функция.}
    \textbf{Корреляционная функция} (КФ) --- отношение измеренных двухчастичного инклюзивного и одночастичного инклюзивного спектров.
    \[C = \frac{A}{B} = \frac{A_{wei}}{A}\]
    \[A_{wei} = femtoweight \cdot A\]
    \[femtoweight = 1 + \cos\frac{(\Delta P\cdot\Delta r)}{hc}\]
    \centering, где $\Delta$ P --- разница 4-импульсов, а $\Delta$ r --- разница 4-векторов
\end{frame}

\begin{frame}{Введение. 1-мерная КФ.}
    В случае 1-мерной КФ используется инвариантный импульс пары:
    \[q_{inv} = \sqrt{{(P_{a} - P_{b})}^{2} - {(E_{a} - E_{b})}^{2}}\]
    Тогда КФ имеет следующий вид:
    \[C(q_{inv}) = \frac{A(q_{inv})}{B(q_{inv})} = \frac{A_{wei}(q_{inv})}{A(q_{inv})}\]
    А аппроксимирующая её кривая:
    \[C(q_{inv}) = 1 + \lambda\exp{(-\frac{R_{inv}^{2} q_{inv}^{2}}{hc})}\]
\end{frame}

\begin{frame}{Введение. 3-мерная КФ.}
    В случае 3-мерной КФ используется:
    \[q = p_{a} - p_{b}\]
    \[P = p_{a} + p_{b}\]
    И система координат LCMS в которой любой 4-вектор $V$ из СЦМ представим в виде:
    
    \[V_{out } = (P_x V_x+P_y V_y) / P_T\]
    \[V_{side } = (P_x V_y-P_y V_x) / P_T\]
    \[V_{long} = (P_0 V_z-P_z V_0) / M_T\]
    \[\text{где } M_T^2=P_0^2-P_z^2 \text{ \;,a } P_T^2=P_x^2+P_y^2\]
    
\end{frame}

\begin{frame}{Введение. 3-мерная КФ.}
    Тогда 3-мерная КФ имеет вид:
    \[
        C = 
        \frac{
            A(q_{out}, q_{side}, q_{long})
        }{
            B(q_{out}, q_{side}, q_{long})
        } = 
        \frac{
            A_{wei}(q_{out}, q_{side}, q_{long})
        }{
            A_(q_{out}, q_{side}, q_{long})
        }
    \]
    А аппроксимирующая её кривая:
    \[
        C(q) = 1 + \lambda \exp{(
            -\sum_{o, s, l}
            \frac{
                R_{ij}q_{i}q_{j}
            }{hc}
        )}
    \]
    В самом простом случае (без учета недиагональных элементов)
    \[
        C(q_{out}, q_{side}, q_{long}) = 
        1 + \lambda \exp{(
            -\frac{
                R_{out}^{2} q_{out}^{2} - 
                R_{side}^{2} q_{side}^{2} - 
                R_{long}^{2} q_{long}^{2}
            }{hc}
        )}
    \]
\end{frame}


\begin{frame}{Настройка модели и отбор частиц.}
	Для генерации была использованна транспортная модель \"SMASH\" со следующими настройками:
    \begin{itemize}
        \item Столкновение: 
            $^{197}_{79}\!Au + ^{197}_{79}\!Au$
        \item При энергии центра масс: 
            $\sqrt{s_{NN}} = 3$ GeV
    \end{itemize}
    А для построения КФ смешивались частицы, удовлетворяющие следующим требованиям:
    \begin{itemize}
        \item Пионы (+) и (-). 
        \item Произошло кинетическое вымораживание.
        \item $P_{total}$ $\in$ [0.15, 1.5] GeV/c
        \item $P_{T}$ $\in$ [0.15, 1.5] GeV/c
        \item $\eta$ $\in$ [-1.5, 1.5]
    \end{itemize}
    Также мы не учитываем события с упругим взаимодействием.
\end{frame}

\begin{frame}{Гистограммы событий.Прицельный параметр}
    \begin{figure}[H]
        \includegraphics[width=0.7\textwidth]{fig/impact_parameter.png}
        \caption{Распределене событий по прицельному параметру}\label{fig:sample}
    \end{figure}
\end{frame}

\setcounter{figure}{3}
\begin{frame}{Гистограммы частиц. 1.Координаты}
    \begin{figure}[H]
        \centering
        \begin{minipage}[b]{0.4\textwidth}
            \centering
            \includegraphics[width=\textwidth]{fig/x.png}
        \end{minipage}
        \begin{minipage}[b]{0.4\textwidth}
            \centering
            \includegraphics[width=\textwidth]{fig/y.png}
        \end{minipage}
        \begin{minipage}[b]{0.4\textwidth}
            \centering
            \includegraphics[width=\textwidth]{fig/z.png}
        \end{minipage}\label{fig:coordinates}
    \end{figure}
\end{frame}

\begin{frame}{Гистограммы частиц. 2.Время кинетического вымораживания}
    \begin{figure}[H]
        \includegraphics[width=0.7\textwidth]{fig/t.png}\label{fig:kineticfreezouttime}
        \caption{Распределение по времени наступления кинетического вымораживания}
    \end{figure}
\end{frame}

\setcounter{figure}{6}
\begin{frame}{Гистограммы частиц. 3.Импульсы}
    \begin{figure}[H]
        \centering
        \begin{minipage}[b]{0.4\textwidth}
            \centering
            \includegraphics[width=\textwidth]{fig/px.png}
        \end{minipage}
        \hfill
        \begin{minipage}[b]{0.4\textwidth}
            \centering
            \includegraphics[width=\textwidth]{fig/py.png}
        \end{minipage}
        \hfill
        \begin{minipage}[b]{0.4\textwidth}
            \centering
            \includegraphics[width=\textwidth]{fig/pz.png}
        \end{minipage}\label{fig:momentums}
    \end{figure}
\end{frame}

\begin{frame}{Гистограммы частиц. 4.Полный и поперечный импульсы}
    \begin{figure}[H]
        \centering
        \begin{minipage}[b]{0.47\textwidth}
            \centering
            \includegraphics[width=\textwidth]{fig/p_total.png}
        \end{minipage}
        \begin{minipage}[b]{0.47\textwidth}
            \centering
            \includegraphics[width=\textwidth]{fig/pt.png}
        \end{minipage}\label{fig:total_and_transvmomentums}
        \caption{Распределение по полному и поперечному импульсу для пионов}
    \end{figure}
\end{frame}

\begin{frame}{Гистограммы частиц. 5. Быстрота и псевдобыстрота}
    \begin{figure}[H]
        \centering
        \begin{minipage}[b]{0.47\textwidth}
            \centering
            \includegraphics[width=\textwidth]{fig/rapidity.png}
        \end{minipage}
        \begin{minipage}[b]{0.47\textwidth}
            \centering
            \includegraphics[width=\textwidth]{fig/pseudorapidity.png}
        \end{minipage}\label{fig:pseudo_and_rapidity}
        \caption{Распределение по быстроте и псевдобыстроте для пионов}
    \end{figure}
\end{frame}



\begin{frame}{Сравнение вычисленных параметров при разных центральностях}
    \begin{figure}[H]
        \centering
        \begin{minipage}[b]{0.47\textwidth}
            \centering
            \includegraphics[width=\textwidth]{fig/c_Rinv_vs_kt.png}
        \end{minipage}
        \begin{minipage}[b]{0.47\textwidth}
            \centering
            \includegraphics[width=\textwidth]{fig/c_lambda_vs_kt.png}
        \end{minipage}\label{fig:Rinv_lambda_4}
    \end{figure}

    $R_{inv}$ и $\lambda$ уменьшается с $k_{t}$ и от центральных к периферическим столкновениям
\end{frame}

\begin{frame}{1-мерная проекция 3-мерной КФ}
    \begin{figure}[H]
        \includegraphics[width=0.7\textwidth]{fig/c_1DCF_min.png}\label{fig:1DCF}
    \end{figure}
\end{frame}

\begin{frame}{2-мерная проекция 3-мерной КФ}
    \begin{figure}[H]
        \includegraphics[width=0.7\textwidth]{fig/c_2DCF_min.png}\label{fig:1DCF}
    \end{figure}
\end{frame}

\begin{frame}{Сравнение вычисленных параметров из 3-мерной КФ при разны центральностях}
    \begin{figure}[H]
        \includegraphics[width=0.65\textwidth]{fig/c_out_side_long_lambda.png}\label{fig:1DCF}
    \end{figure}
\end{frame}

\begin{frame}{Заключение}
    \begin{itemize}
	\item Был сгенерирован набор частиц в транспортной модели “SMASH”.
        \item Были построены и изучены гистограммы для базового анализа столкновения.
        \item На базе сгенерированных модельных данных был освоен метод фемтоскопических корреляций.
        \item Были построены корреляционные функции для одномерного и трехмерного анализа.
        \item Были получены зависимости радиусов и силы корреляции от поперечного импульса и центральности.
    \end{itemize}
\end{frame}

\begin{frame}{Цели}
    \begin{itemize}
	\item  Набрать более полную статистику для высокой центральности и для разных значений $\sqrt{S_{NN}}$.
        \item Изучить асимметрию. В частности для направления out. 
        \item Проверить соответствие модельных и экспериментальных данных.
    \end{itemize}
\end{frame}

\begin{frame}
	\Large \textbf{Спасибо за внимание!}
\end{frame}

\begin{frame}{Дополнительные слайды. КФ для всех центральностей и $k_{t}$ классов}
    \begin{figure}[H]
        \includegraphics[width=0.65\textwidth]{fig/c_all_1d_CF.png}\label{fig:1DCF}
    \end{figure}
\end{frame}

\begin{frame}{Дополнительные слайды. 3-мерные КФ для всех центральностей и $k_{t}$ классов}
    \begin{figure}[H]
        \includegraphics[width=0.65\textwidth]{fig/c_all_3d_CF.png}\label{fig:1DCF}
    \end{figure}
\end{frame}

\begin{frame}{Дополнительные слайды. 1-мерные проекции 3-мерных КФ для всех центральностей и $k_{t}$ классов, на ось out}
    \begin{figure}[H]
        \includegraphics[width=0.65\textwidth]{fig/c_all_3d_CF_proj_out.png}\label{fig:1DCF}
    \end{figure}
\end{frame}

\begin{frame}{Дополнительные слайды. 1-мерные проекции 3-мерных КФ для всех центральностей и $k_{t}$ классов, на ось side}
    \begin{figure}[H]
        \includegraphics[width=0.65\textwidth]{fig/c_all_3d_CF_proj_side.png}\label{fig:1DCF}
    \end{figure}
\end{frame}

\begin{frame}{Дополнительные слайды. 1-мерные проекции 3-мерных КФ для всех центральностей и $k_{t}$ классов, на ось long}
    \begin{figure}[H]
        \includegraphics[width=0.65\textwidth]{fig/c_all_3d_CF_proj_long.png}\label{fig:1DCF}
    \end{figure}
\end{frame}

\begin{frame}{Дополнительные слайды. 2-мерные проекции 3-мерных КФ для всех центральностей и $k_{t}$ классов, в осях out-side}
    \begin{figure}[H]
        \includegraphics[width=0.65\textwidth]{fig/c_2DCF_out_vs_side.png}\label{fig:1DCF}
    \end{figure}
\end{frame}

\begin{frame}{Дополнительные слайды. 2-мерные проекции 3-мерных КФ для всех центральностей и $k_{t}$ классов, в осях out-long}
    \begin{figure}[H]
        \includegraphics[width=0.65\textwidth]{fig/c_2DCF_out_vs_long.png}\label{fig:1DCF}
    \end{figure}
\end{frame}

\begin{frame}{Дополнительные слайды. 2-мерные проекции  3-мерных КФ для всех центральностей и $k_{t}$ классов, в осях side-long}
    \begin{figure}[H]
        \includegraphics[width=0.65\textwidth]{fig/c_2DCF_side_vs_long.png}\label{fig:1DCF}
    \end{figure}
\end{frame}


